\section{Two-sided marked correlations}\label{sec:correlations}
\label{cor:section}

We now allow marks on both endpoints of the shift.  The integer $J$ used
in this section is sufficiently large and \emph{fixed} as $x\to\infty$;
it is different from the growing determinant parameter in the preceding
sections.  Likewise the active damping parameter $q_0$ below need not
equal the parameter $q$ in the fixed-band weight $\mathcal W$.

Retain the fixed $K$ bands defining $\mathcal W$.  Independently choose
pairwise disjoint active groups $\mathcal P_g$, indexed by
$\mathcal S\sqcup\mathcal B$, such that, for fixed positive constants
$c_0,C_0,v_-,v_+$,
\begin{gather}
 c_0\log L\leq |\mathcal S|,|\mathcal B|\leq C_0\log L,
 \qquad v_-\leq V_g:=\sum_{p\in\mathcal P_g}p^{-1}\leq v_+,
 \label{cor:groups}\\
 \mathcal P_g\subset[\exp(L^{.27}),\exp(L^{.36})]
       \quad(g\in\mathcal S),\qquad
 \mathcal P_g\subset[\exp(L^{.39}),\exp(L^{.46})]
       \quad(g\in\mathcal B).\notag
\end{gather}
Every big group is the set of all primes in an interval.  Write
$\mathcal P=\bigcup_g\mathcal P_g$, $s_P=|\mathcal S|+|\mathcal B|$,
and define
\[
 n_*:=\prod_{p\notin\mathcal P}p^{v_p(n)},\qquad
 \omega_g(n):=\sum_{p\in\mathcal P_g}\1_{p\mid n},\qquad
 \omega_P(n):=\sum_g\omega_g(n),\qquad
 \mu_g(p):=\frac1{V_gp}.
\]
Thus $n_*$ removes the entire active part, including its multiplicities.
For a vector $\mathbf t=(t_g)$ of nonnegative integers, a represented
list at $n$ consists of $t_g$ ordered distinct prime divisors of $n$
from each group.  For a function $\mathbf b$ of these lists put
\begin{equation}\label{cor:marked-weight}
 W_{\mathbf t}^{\mathbf b}(n)
  :=q_0^{\omega_P(n)-\sum_g t_g}
       \prod_gV_g^{-t_g}
       \sum_{\text{represented lists at }n}\mathbf b(\mathbf p),
 \qquad 0<q_0<1.
\end{equation}
The weight is zero if no list exists.  Write $W_{\mathbf t}$ when
$\mathbf b=1$, and $W_1=W_{(1,\ldots,1)}$.  For every fixed $m'$,
\begin{equation}\label{cor:marked-bound}
 |W_{\mathbf t}^{\mathbf b}(n)|\leq W_{\mathbf t}(n)
 \leq L^{C(m')}
 \quad\bigl(t_g\leq m',\ |\mathbf b|\leq1\bigr).
\end{equation}
Write $(u)_t=u(u-1)\cdots(u-t+1)$, with $(u)_0=1$.
Indeed, each factor
$q_0^{u-t}(u)_t/V_g^t$, for $u\geq t$, is bounded by a constant
depending only on $m',q_0,v_-$, and there are $O(\log L)$ factors.

\begin{theorem}[Two-sided correlation]\label{thm:two-sided}
Fix $d\geq1$, $\eps>0$, and intervals $I_1,\ldots,I_d$ whose lower
endpoints are at least $x^\eps$ and whose upper endpoints have product
at most $x^{1-\eps}$.  Let $\mathbf l=(l_1,\ldots,l_d)$ run over
ordered tuples of distinct primes with $l_i\in I_i$, and put
\begin{equation}\label{cor:centered-core}
 l:=l_1\cdots l_d,\qquad
 C_x:=\sum_{\mathbf l}\frac1l,\qquad
 F_x(n):=\mathcal W(n)
       \left(\sum_{\mathbf l}\1_{l\mid n}-C_x\right).
\end{equation}
Suppose $G(n)=G(n_*)$ and
$|G(n)|\leq L^C\tau(n_*)^C$, for fixed $C$.  For every fixed smooth
compactly supported $\Psi:(0,\infty)\to\C$ and every fixed $A>0$,
\begin{equation}\label{eq:two-sided}
 \sum_{w\geq1}\frac{\Psi(w/x)}w
       F_x(w)G(w+1)W_1(w)W_1(w+1)\ll_A L^{-A}.
\end{equation}
The constants and sufficiently-large-$x$ threshold may depend on all
the fixed data, including the fixed inert bands, but the estimate is
uniform over the active groups and slot intervals satisfying the stated
bounds.  The slots have no joint restriction other than distinctness.
\end{theorem}

The harmonic prime estimates give $C_x=O_\eps(1)$.  On every fixed
positive-power range for $n$, there are only $O_\eps(1)$ prime divisors
in the slot ranges.  Since $\mathcal W$ is bounded for fixed $K,q$,
$F_x$ is bounded there.  Also $F_x(pn)=F_x(n)$ for every active prime
$p$: the active groups are disjoint both from the inert bands and from
the long-prime slots.  These two facts will give the stronger endpoint
hypotheses required by the graph argument.

The proof has two stages.  First, the graph inputs from S give
arbitrary logarithmic savings for a signed combination of correlations.
After common prime labels have been divided out of both endpoints,
one term, called the raw term, is the unit-shift correlation
in \eqref{eq:two-sided}; the other terms have shifts that are products
of independently sampled big-group primes.  We call these other terms
the comparison correlations.  We then bound each comparison separately,
using the centered first endpoint, and recover the unit-shift term by
subtraction.  We state the graph inputs with their exact measures before
making this reduction.  The new endpoint estimates begin once the
comparison correlations have been identified.

\subsection{The exact graph inputs}

For clarity we state the general graph contracts, rather than using a
correlation theorem with more restrictive endpoint hypotheses.  In this
subsection only, replace $.27,.36,.39,.46$ in
\eqref{cor:groups} by arbitrary fixed $0<a<b<c<d_0<.47$, and fix an
integer $\ell\geq1$.  Extend the divisibility definitions to all
$n\in\Z$, with every group prime dividing zero.  Put $M=J+\ell$.
A pattern $\nu$ chooses sets
\[
 C_g\subset\{1,\ldots,J\},\qquad C_g=\varnothing\ (g\in\mathcal S),
 \qquad |C_g|\leq m_0\ (g\in\mathcal B),
 \qquad t_g=\ell+|C_g|.
\]
In each of the first $J$ slots outside $C_g$, the source and target
lists have the same prime label.  The product of these shared labels
is $D_R$.  In each slot of $C_g$ take an independent
auxiliary label of law $\mu_g$; their product is $D_C$.  The dilation
$D=D_RD_C$ therefore contains $J$ labels from each group, counted with
multiplicity.  A coefficient $K_\nu$ may depend only on the ordered
big-group source lists, target lists, and auxiliary free labels.

The physical Hilbert space consists of states $(n,\mathbf p)$ with
$n\in\Z$ and $M$ ordered distinct divisors from every group.  Each
state has mass $\prod_gV_g^{-M}$, with counting measure in $n$.
Fix an integer $k$ with $0<|k|\leq L^{C_2}$.  A row operation samples the auxiliary labels,
sets $n'=n+kD$, and retains the shared labels while summing over
physical target lists with coefficient $\prod_gV_g^{-t_g}$.
Every auxiliary free label is forbidden from both endpoint lists;
auxiliary labels may coincide with one another.  The row multiplier is
\begin{equation}\label{cor:physical-multiplier}
 K_\nu D^{i\zeta}
 q_0^{(\omega_P(n)-Ms_P)/2}
 q_0^{(\omega_P(n')-Ms_P)/2}.
\end{equation}
Add the patterns and average independent permutations of the $M$ slots
in every group at both ends.  This defines $\mathcal A_\zeta$.
In particular, the joint measure of the two lists in group $g$ is
$V_g^{-M-t_g}$ in either direction.  A row operator integrates the
target function and returns a function at the source.

The ideal Hilbert space is
\[
 \mathcal H_{\rm id}
  =L^2\left(\prod_{g\in\mathcal B}\mathcal P_g^M,
                \bigotimes_{g\in\mathcal B}\mu_g^{\otimes M}\right).
\]
Repetitions are allowed here.  Retain the shared coordinates and sample
the unshared target and auxiliary coordinates independently.  If
$D_{\mathcal B}$ is the big-group part of $D$, use multiplier
$K_\nu D_{\mathcal B}^{i\zeta}\e(\Theta D_{\mathcal B})$.
The sum, symmetrized at both ends, is $\mathcal T_\zeta(\Theta)$.

\begin{proposition}[General graph transference input]
\label{prop:graph-transfer}
Assume $\sum_\nu\sup|K_\nu|\leq L^{C_1}$, where $m_0,C_1$ are
fixed independently of $J$.  For every fixed $E>0$ there is
$E_{\rm id}>0$, depending only on $E$ and the fixed group data,
such that
\[
 \sup_{\Theta\in\R}
    \|\mathcal T_\zeta(\Theta)\|_{2\to2}\leq L^{-E_{\rm id}}
\]
has the following consequence for all sufficiently large fixed $J$.
Let $I=[X,2X)$, with $x^\gamma\leq X\leq x^{1/\gamma}$ for fixed
$\gamma>0$.  If $f$ is supported on positions in $I$, is independent
of the chosen marks, and satisfies
\[
 \sup|f|\leq\exp(O(\sqrt L)),\qquad
 \|f\|,\|g\|\leq X^{1/2}L^{C_3},
\]
then
\begin{equation}\label{cor:transferred-pairing}
 |\langle f,\mathcal A_\zeta g\rangle|
       \ll XL^{-E+2C_3}.
\end{equation}
The lower bound on $J$ may also depend on $m_0,C_1$.  The threshold
for $x$ may depend on $J,C_2,\gamma$.  No further restriction on
$\zeta$ is imposed beyond the ideal norm hypothesis.  Taking absolute
values of the individual transition coefficients gives row and column
sums at most $L^{C_{\rm abs}}$, where $C_{\rm abs}$ is independent
of $J$.
\end{proposition}

This is the combination of the local transference theorem, endpoint
pairing corollary, and physical Schur bound of
S~\cite[Theorem~3.5, Corollary~3.11, and Lemma~3.4]{SmoothShifted}.
Its parameter $q$ is our $q_0$; its laws, physical state masses,
two-sided symmetrization, and ban on free endpoint labels are exactly
those defined above. In the absolute bound,
if a target has $y\geq M$ divisors in a group, its unshared list sum
is bounded by
\[
 \sup_{z\geq0}V_g^{-t_g}(z+t_g)_{t_g}q_0^{z/2}=O_{m_0,\ell}(1).
\]
The source damping is at most one.  Products over $O(\log L)$ groups
and the pattern mass give the asserted exponent independent of $J$;
the symmetric joint normalization gives the column bound as well.

\begin{proposition}[Residual ideal family input]\label{prop:ideal-family}
For every $E_{\rm id}>0$ and fixed $C_4\geq0$, there are fixed
$m_0,J_0,C_1$ such that, for each fixed $J\geq J_0$, a family
consisting of the raw pattern $C_g=\varnothing$, $K_0=1$, and signed
comparison patterns satisfies
\begin{equation}\label{cor:ideal-bound}
 \sum_\nu\sup|K_\nu|\leq L^{C_1},\qquad
 \sup_{\substack{\Theta\in\R\\|\zeta|\leq L^{C_4}}}
       \|\mathcal T_\zeta(\Theta)\|_{2\to2}\leq L^{-E_{\rm id}}.
\end{equation}
The family is independent of $\Theta,\zeta$; its defining parameters
are independent of $J$.

Split the big groups into two blocks.  Every comparison frees nonempty
sets of probes $A,B$ in the respective blocks, with coefficient
\begin{equation}\label{cor:comparison-kernel}
 -(-1)^{|A|+|B|}\mathcal K_A(D_A,X_A)\mathcal K_B(D_B,X_B).
\end{equation}
Here $D_A,D_B$ are the free products, $X_A$ is the corresponding
target-mark product, and $X_B$ the corresponding source-mark product.
No shared or final $\ell$ label enters these kernels.  For a nonempty
slot set $\Lambda$ with at most $m_0$ slots in each group,
\begin{equation}\label{cor:cell-kernel}
 \mathcal K_\Lambda(D,X)
 =\sum_{j:\nu_{\Lambda j}\geq\xi}
   \frac{\1_{\log D\in I_j}\1_{\log X\in I_j}}{\nu_{\Lambda j}}
   \sum_{\substack{\chi\text{ primitive}\\\cond(\chi)\leq Q}}
        \overline{\chi(D)}\chi(X),
 \qquad I_j=[jh,(j+1)h).
\end{equation}
The number $\nu_{\Lambda j}$ is the cell probability for the
independent slot laws.  The parameters $Q,h^{-1},\xi^{-1}$ are fixed
powers of $L$, with $h\leq1$.  There are at most $Q^2$ characters
and $O_{m_0}(1+L^{d_0}\log L/h)$ nonempty cells, and
\[
 \sup|\mathcal K_\Lambda|\leq Q^2\xi^{-1},\qquad
 \sup_X\E_D|\mathcal K_\Lambda(D,X)|\leq Q^2.
\]
More precisely, for any prescribed fixed $A_0>0$, the parameters may
be chosen, independently of $J$, so that for arbitrary $L^2$ tuple
tests $f(X_B),g(X_A)$ and all $\theta\in\R$, $|\zeta|\leq L^{C_4}$,
\begin{align}
 \big|\E\,\overline{f(X_B)}g(X_A)
 \big[&(X_AX_B)^{i\zeta}\e(\theta X_AX_B)\notag\\
 &-\mathcal K_A(D_A,X_A)\mathcal K_B(D_B,X_B)
       (D_AD_B)^{i\zeta}\e(\theta D_AD_B)\big]\big|
 \leq L^{-A_0}\|f\|_2\|g\|_2.
 \label{cor:comparison-contract}
\end{align}
The four label tuples in this expectation are independent; the tests
need not depend only on their products.  Expanding the two kernels,
including all patterns, has total coefficient mass and number of terms
bounded by fixed powers of $L$.
\end{proposition}

This is S~\cite[Theorem~4.1 and Lemma~4.3]{SmoothShifted}.
The kernels
in \eqref{cor:cell-kernel} are two global cell and character expansions;
there is not a separate character expansion for every group.  We use
these two imported propositions with $\ell=k=1$ and the exponents in
\eqref{cor:groups}.

\subsection{Endpoint bounds and removal of shared labels}

On a physical state $\omega_P(n)\geq Ms_P$.  Use endpoint vectors
with values $\overline{F_x(n)}$ and $G(n')$, respectively, and with
one extra half-damping factor at each endpoint.  The conjugation
compensates for that in the Hilbert-space pairing.  These extra factors
are at most one and are independent of the lists.  For a fixed ordered
physical list with squarefree product $P$, $n_*\mid n/P$.  The divisor
moment bound of Lemma~\ref{lem:moment-bounds} gives
\[
 \sum_{\substack{n\asymp X\\P\mid n}}|G(n)|^2
 \leq L^{2C}\sum_{v\ll X/P}\tau(v)^{2C}
 \ll (X/P)L^{C_*}.
\]
The same statement holds for $F_x$, which is bounded.  Here
$P\leq\exp(O_J(L^{.46}\log L))=x^{o(1)}$, and the normalized list
sum satisfies
\begin{equation}\label{cor:list-norm}
 \prod_gV_g^{-M}\sum_{\text{physical lists}}\frac1P
 \leq\prod_g\left(\sum_{p\in\mathcal P_g}\frac1{V_gp}\right)^M=1.
\end{equation}
Consequently both endpoint norms on an enlarged dyad are at most
$X^{1/2}L^{C_3}$, where $C_3$ is independent of $J$ and $m_0$.
The first vector has bounded supremum.  It is essential here to keep
$q_0^{(\omega_P(n)-Ms_P)/2}\leq1$ intact; separating its two powers
would introduce an unnecessary exponent depending on $J$.

Insert $\Psi(n/(xD))/n$ into the physical edge pairing.  Its support
has $n\asymp n'\asymp xD=x^{1+o(1)}$.  Partition $n$ into $O(L)$
dyads $[X,2X)$, restricting $n'$ to fixed enlargements.  Let
$\phi(u)=\Psi(e^u)$ and use the convention
$\widehat\phi(\zeta)=\int\phi(u)e^{-i\zeta u}\,\dd u$.  Then
\begin{equation}\label{cor:insertion-separation}
 \Psi\left(\frac n{xD}\right)
   =\frac1{2\pi}\int_{\R}\widehat\phi(\zeta)
       (n/x)^{i\zeta}D^{-i\zeta}\,\dd\zeta.
\end{equation}
Put $X/n$ into the first vector and $1/X$ outside.  The endpoint
bounds just proved are unchanged.  Fix the cutoff $|\zeta|\leq L$
now, so $C_4=1$ suffices in Proposition~\ref{prop:ideal-family}.

The choices have the following order.  After the desired saving $A$
and the original data, fix $C_3$, then a transfer saving $E$ large
enough to cover $A$, $2C_3$, and the fixed dyadic costs.  Choose
$E_{\rm id}$ from Proposition~\ref{prop:graph-transfer}, obtain
$m_0,J_0,C_1$ from Proposition~\ref{prop:ideal-family}, and finally
choose a fixed $J$ satisfying both lower bounds.  For the tail of
\eqref{cor:insertion-separation}, the absolute Schur estimate and
the endpoint norms cost at most a fixed power of $L$.  Since
\[
 \int_{|\zeta|>L}|\widehat\phi(\zeta)|\,\dd\zeta
       \ll_N L^{-N}
\]
for every fixed $N$, its differentiation order can be chosen after
the family is fixed.  The cutoff exponent does not have to change.
Thus the total residual edge pairing is $O_A(L^{-A})$, with as much
extra fixed saving as will be needed below.

For a single pattern put $n=D_Rw$, $n'=D_Rw'$.  The shift becomes
$w'=w+D_C$.  Away from shared-label collisions,
\begin{equation}\label{cor:stripping-normalization}
 \omega_g(D_Rw)-M=\omega_g(w)-t_g,
 \qquad
 V_g^{-M-t_g}=V_g^{-(M-t_g)}V_g^{-2t_g},
 \qquad \frac1n=\frac1{D_Rw}.
\end{equation}
The factor $D_R^{-1}$ changes each shared sum into its probability
law $\mu_g$, and the remaining list factors and damping are exactly
$W_{\mathbf t}(w)W_{\mathbf t}(w')$.  The endpoint cores are invariant
under these labels.  The comparison coefficients involve no shared
labels.  Consequently the stripped expression is
\begin{equation}\label{cor:stripped}
 \E_{D_C}\sum_{w\geq1}\frac{\Psi(w/(xD_C))}{w}
       F_x(w)G(w+D_C)
       W_{\mathbf t}(w)W_{\mathbf t}(w+D_C)
       \E_{\mathbf p(w),\mathbf p(w+D_C)}K_\nu,
\end{equation}
where the last expectation is uniform on each endpoint's ordered
unshared lists.  The independent permutations have total mass one and
only rename the ordered slots, so they introduce no factorial factor.

We justify restoring all the restrictions suppressed in this formula.
There are $O_J(\log L)$ shared or free slots.  Every group atom
has size $O(\exp(-L^{.27}))$.  Conditional on $w,w'$ and the free
labels, the harmonic mass of a shared/shared coincidence, a
shared/free coincidence, or a shared prime dividing $ww'$ is therefore
\begin{equation}\label{cor:shared-collision}
 O_J\bigl(L^{C_J}\exp(-L^{.27})\bigr).
\end{equation}
For the last assertion use that an integer of size $x^{O(1)}$ has at
most $O(L)$ distinct prime factors.  On actual overlap exceptions the
unshared labels still divide $w,w'$.  Changing the damping loses at
most $q_0^{-2Ms_P}=L^{O(J)}$.  The marked bounds and Cauchy with
fixed divisor moments on the translated dyads bound the remaining
harmonic sum by a fixed log power.  Thus
\eqref{cor:shared-collision} also controls these actual exceptions.

If a free prime occurs in either unshared endpoint list, it divides
$D_C$ and one of $w,w+D_C$, hence both.  It is at least
$\exp(L^{.39})$.  On $w\asymp xD_C$ the count of its multiples is
$O(xD_C/p)$; Cauchy and a fixed higher divisor moment bound the
corresponding weighted harmonic mass by $L^{O(1)}p^{-1/2}$.
Summing over the $O_{m_0}(\log L)$ free slots costs
$L^{O(1)}\exp(-cL^{.39})$.  This also restores any coincidence between
the two unshared endpoint lists, since such a label must divide their
difference $D_C$.  All errors remain smaller than every negative log
power after the fixed pattern cost.  They are incurred after $J$ is
fixed and do not change the exponent $C_3$ used in transference.
For the raw pattern, $D_C=1$ and $t_g=1$, so
\eqref{cor:stripped} is precisely the left side of
\eqref{eq:two-sided}.

The transferred estimate therefore controls the desired correlation
plus the signed comparison correlations \eqref{cor:stripped}.  It
remains to bound the latter separately.  Their shifts contain two
nonempty products of free big-group labels; these products will give
a bilinear Fourier multiplier.  All cancellation required at an
endpoint will come from the centered function $F_x$.

\subsection{Comparison multipliers and local energy}

Expand the two kernels in \eqref{cor:comparison-kernel}.  Their
contract supplies only a fixed log-power total cost.  In each term the
two nonempty free products are independent and lie in log cells of
width at most one, say $D_A\asymp U_1$, $D_B\asymp U_2$.  Put
\begin{equation}\label{cor:comparison-scales}
 H'=U_1U_2,\qquad Y=xH',\qquad
 U_i\geq\exp(L^{.39}-O(1)),\qquad
 H'\leq\exp(O_{m_0}(L^{.46}\log L)).
\end{equation}
The cell and character expansion factors into separate bounded tests
on the two free tuples and on each endpoint's big marks.  Insert fixed
smooth cutoffs at scale $Y$ at both endpoints.  Fourier inversion in
logarithmic size separates $\Psi(w/(xD_AD_B))/w$ as in
\eqref{cor:insertion-separation}.  After a factor $1/Y$, it is enough
to estimate combinations of
\begin{equation}\label{cor:comparison-sum}
 \E\,b_1(D_A)b_2(D_B)
       \sum_w\mathcal U(w)\mathcal V(w+D_AD_B),
 \qquad |b_i|\leq1.
\end{equation}
Tuple tests may be used in place of product tests; conditioning on the
product produces the notation here.  The endpoint sequences have
smooth cutoffs at $Y$, fixed log-power twists, and the form of their
respective invariant cores times $W_{\mathbf t}^{\mathbf b}$, with
$|\mathbf b|\leq1$ depending only on big marks.  Fixed divisor moments
give
\begin{equation}\label{cor:endpoint-squares}
 \sum_w|\mathcal U(w)|^2+\sum_w|\mathcal V(w)|^2
       \ll YL^{C_5}.
\end{equation}
These estimates also bound the separated Fourier tails absolutely:
on any translated time interval, the corresponding collapsed
Dirichlet polynomials obey the coefficient mean-square estimate.
Smooth Fourier decay can therefore give any required precision.
All exponents chosen from now on may depend on the fixed ideal family.

The additive multiplier for \eqref{cor:comparison-sum} is
\[
 \mathfrak m(\theta)
   :=\E\,b_1(D_A)b_2(D_B)\e(\theta D_AD_B),
 \qquad |\mathfrak m(\theta)|\leq1.
\]
The product probability of a tuple using $k_g\leq m_0$ slots per
group is, by unique factorization,
\[
 \rho(u)=\frac1u\prod_g
       \frac{k_g!}{V_g^{k_g}\prod_p a_p!}\leq\frac{L^{C_7}}u.
\]
Since its total mass is at most one, each collapsed sequence on
$u\asymp U_i$ has squared coefficient norm at most $L^{C_7}/U_i$.
The elementary bilinear Fourier estimate (Cauchy followed by the
spacing estimate for $\|\theta u\|^{-1}$) consequently gives, if
$(h,r)=1$ and $|\theta-h/r|\leq r^{-2}$,
\begin{equation}\label{cor:bilinear-multiplier}
 |\mathfrak m(\theta)|
  \ll L^{C_7}
    \left(\frac{(U_1/r+1)(U_2+r\log(2r))}{H'}\right)^{1/2}.
\end{equation}
For completeness, Cauchy in the variable on the interval of length
$O(U_2)$, expansion, and the geometric-sum bound give an unnormalized
square at most
\[
 \|a\|_2^2\|b\|_2^2
 \left(U_2+\sum_{1\leq v\ll U_1}
                   \min\{U_2,\|\theta v\|_{\R/\Z}^{-1}\}\right).
\]
For $r\geq2$, split the $v$-range into $O(1+U_1/r)$ blocks of
diameter at most $r/2$.  Distinct points in one block have residues
$\theta v$ separated by at least $1/(2r)$, because
$|\theta-h/r|\leq r^{-2}$.  The sum in a block is at most
$O(U_2+r\log(2r))$.  The case $r=1$ is immediate from the trivial
bound $U_2$ on each summand.  This proves
\eqref{cor:bilinear-multiplier} with the two coefficient norms above.

Given any required saving, choose $C_6$ sufficiently large and set
\begin{equation}\label{cor:major-arcs}
 H=H'/L^{C_6},\qquad
 \mathfrak M_H=\bigcup_{\substack{k\leq L^{C_6}\\(h,k)=1}}
         \{\theta:|\theta-h/k|\leq H^{-1}\}.
\end{equation}
Dirichlet approximation with denominator bound
$\lfloor H'/L^{C_6-1}\rfloor$ gives a fraction $h/r$ with
$|\theta-h/r|\leq r^{-2}$.  If $r\leq L^{C_6}$, this places
$\theta$ in $\mathfrak M_H$.  Otherwise
$L^{C_6}<r\ll H'/L^{C_6-1}$, and the expression under the square
root in \eqref{cor:bilinear-multiplier} is bounded by
\[
 \frac1r+\frac{\log(2r)}{U_2}
       +\frac1{U_1}+\frac{r\log(2r)}{H'}.
\]
The middle terms save every log power by
\eqref{cor:comparison-scales}, and the other terms give any desired
fixed saving by increasing $C_6$.  Thus $\mathfrak m$ is arbitrarily
log-power small off $\mathfrak M_H$.

Use the Fourier convention
$\widehat a(\theta)=\sum_w a_w\e(\theta w)$.
Parseval and \eqref{cor:endpoint-squares} handle the complement of
$\mathfrak M_H$.  On its $O(L^{2C_6})$ arcs, Cauchy and the global
energy of $\mathcal V$ reduce the problem to
\begin{equation}\label{cor:arc-energy-target}
 \int_{|\beta|\leq H^{-1}}
     |\widehat{\mathcal U}(h/k+\beta)|^2\,\dd\beta
       \ll_{A'}YL^{-A'}
\end{equation}
for every fixed $A'$.  The scale conversion we need is as follows.

\begin{lemma}[Local Mellin energy]\label{lem:local-mellin}
Let $a_w$ be supported on $w\asymp Y$, where $Y=x^{1+o(1)}$,
and suppose $H=x^{o(1)}\to\infty$.  For every fixed integer $j\geq1$,
\begin{equation}\label{cor:local-mellin}
 \int_{|\beta|\leq H^{-1}}|\widehat a(\beta)|^2\,\dd\beta
 \ll_j\frac1Y\int_{\R}(1+H|t|/Y)^{-j}
       \left|\sum_w a_ww^{it}\right|^2\,\dd t
       +\left(\frac HY\right)^2\sum_w|a_w|^2.
\end{equation}
\end{lemma}

\begin{proof}
Choose a smooth bump $K$ of
integral one, supported close enough to zero that its additive Fourier
transform has modulus at least $1/2$ on $[-1,1]$.  Plancherel gives
\[
 \int_{|\beta|\leq H^{-1}}|\widehat a(\beta)|^2\,\dd\beta
 \ll H^{-2}\int_\R
       \left|\sum_w a_wK((w-v)/H)\right|^2\,\dd v.
\]
Only $v\asymp Y$ contributes.  On the union of the relevant supports,
replace $K((w-v)/H)$ by $K(v\log(w/v)/H)$: their arguments differ
by $O(H/Y)$.  Cauchy over the $O(H)$ available $w$, and integration
over the $O(H)$ centers for each $w$, give normalized squared error
$O((H/Y)^2\sum|a_w|^2)$.

Write $v=Ye^u$, $h_0=H/Y$, and $P_a(t)=\sum_w a_ww^{it}$.
Fourier inversion for the new kernel reads
\[
 \sum_w a_wK(v\log(w/v)/H)
 =\frac{h_0}{2\pi}\int_\R e^{-u}
      \widehat K(h_0e^{-u}t/(2\pi))v^{-it}P_a(t)\,\dd t.
\]
Insert a fixed smooth cutoff in $u$ for $v\asymp Y$.  Two
integrations by parts in $u$ and the Schwartz bounds for $\widehat K$
bound the Fourier transform of this amplitude by
$C_j(1+|s|)^{-2}(1+h_0|t|)^{-j}$.  Expand in $s$, apply Minkowski
and Plancherel in $u$, and use $\dd v\ll Y\,\dd u$.  The outside
normalization is $H^{-2}Yh_0^2=Y^{-1}$.  Increasing the decay order
proves \eqref{cor:local-mellin}.
\end{proof}

Apply the lemma to $a_w=\mathcal U(w)\e(hw/k)$.  The additive error
is negligible by \eqref{cor:endpoint-squares}.  The Dirichlet
mean-square bound in Lemma~\ref{lem:moment-bounds} gives, on every
dyadic time interval of length $R$,
\[
 Y^{-2}\int_{R\leq|t|\leq2R}|P_a(t)|^2\,\dd t
       \ll L^{C_5+1}(1+R/Y).
\]
For $T_0=LY/H$, the tail in \eqref{cor:local-mellin} is thus at most
\[
 YL^{C_5+1}\sum_{r\geq0}(2^rL)^{-j}
                      (1+2^rL/H),
\]
which is as small as required on choosing $j$ last.  It remains to
prove, for every fixed $A'>0$,
\begin{equation}\label{cor:SM}
 \int_{|t|\leq T_0}
 \left|Y^{-1}\sum_w\mathcal U(w)\e(hw/k)w^{it}\right|^2\,\dd t
       \ll_{A'} L^{-A'},
 \qquad T_0=\frac{LY}{H}=xL^{C_6+1}.
\end{equation}
Notice that the first term in \eqref{cor:local-mellin} is $Y$ times
the integral in \eqref{cor:SM}, as required by
\eqref{cor:arc-energy-target}.

We now prove this one endpoint estimate.  At low Mellin frequencies,
a finite divisor model makes the centering in $F_x$ cancel the main
term.  At high frequencies, one small-prime mark supplies a polynomial
that is small outside a sparse set of times.  On that sparse set we
treat the positive tuple sum and its subtracted constant separately:
the tuple sum has a long-prime factor, whereas the constant term
requires cancellation on the divisor progressions of the finite model.

\subsection{A divisor model and low Mellin frequencies}

\begin{lemma}[Truncated divisor model]\label{lem:divisor-model}
Let
$H_{\rm marks}(n)=\mathcal W(n)W_{\mathbf t}^{\mathbf b}(n)$,
where $0\leq t_g\leq m'$ for fixed $m'$ and $|\mathbf b|\leq1$.
In particular, zero values of $t_g$ are allowed; these occur when a
small mark is removed below.
There is a divisor sum
\begin{equation}\label{cor:divisor-model}
 H_0(n)=\sum_{d'\mid n}A_{d'},\qquad
 d'\leq\exp(O(L^{.47})),\qquad
 \sum_{d'}\frac{|A_{d'}|}{d'}\leq L^{C_8},
\end{equation}
whose prime divisors all belong to the active or inert groups, such
that, on every positive integer dyad $n\asymp R$,
\begin{equation}\label{cor:model-error}
 \sum_{n\asymp R}|H_{\rm marks}(n)-H_0(n)|^2
       \ll_N R L^{-N}
\end{equation}
for every fixed $N$.  The assertion is uniform in the bounded tuple
test $\mathbf b$, with no regularity assumption on that test.
\end{lemma}

\begin{proof}
Expand both weights into their represented lists.  The inert list has
one prime per band, and the active list has $t_g$ per group.  For a
fixed represented list, the damping is exactly the product of $q$
or $q_0$ over the remaining distinct group primes that divide $n$.
Expand this product as
\[
 \prod_{p\text{ not represented}}(1-(1-q_p)\1_{p\mid n}),
 \qquad q_p\in\{q,q_0\},
\]
and truncate after $h_0=\lfloor L^{.01}\rfloor$ additional primes.
All represented primes are distinct within their groups, and the
additional primes are disjoint from them.  Each resulting divisor has
at most $O_{m'}(\log L)+K+h_0$ prime factors, all at most
$\exp(L^{.46})$.  This proves the support assertion in
\eqref{cor:divisor-model}.

In the reciprocal coefficient sum the normalized represented lists
have total mass at most one, by their harmonic definitions.  The
additional subset sums are at most
\[
 \prod_{p\text{ in all groups}}(1+(1-q_p)/p)
       \leq\exp(O(\log L))=L^{O(1)}.
\]
This proves the coefficient bound, even after absolute values.

Let $\omega_{\rm tot}$ count hits in all active and inert groups.
For fixed list sizes, their normalized counts are at most
$L^{C_9}2^{\omega_{\rm tot}}$: use
$(u)_t=t!\binom ut\leq t!2^u$ in each group.  The discarded
binomial expansion is at most
$2^{\omega_{\rm tot}}\1_{\omega_{\rm tot}\geq h_0}$.
Hence the absolute pointwise error is bounded by
$L^{C_9}4^{\omega_{\rm tot}}\1_{\omega_{\rm tot}\geq h_0}$.
For any fixed $B>1$, expanding $B^{\omega_{\rm tot}}$ into
squarefree divisors and counting their multiples gives
\begin{equation}\label{cor:omega-moment}
 \sum_{n\asymp R}B^{\omega_{\rm tot}(n)}
 \ll R\prod_{p\text{ in all groups}}(1+(B-1)/p)
 \ll R L^{C(B)}.
\end{equation}
Indeed only divisors at most a constant times $R$ can occur, and each
has $O(R/d)$ multiples in the dyad.  Since
$16^u\1_{u\geq h_0}\leq2^{-h_0}32^u$, the square of the error
has total at most $R L^{O(1)}2^{-h_0}$.  This proves
\eqref{cor:model-error} for every fixed $N$.
\end{proof}

Write the first endpoint as
\begin{equation}\label{cor:U-form}
 \mathcal U(w)=\psi_1(w/Y)w^{i\sigma}
       \left(\sum_{\mathbf l}\1_{l\mid w}-C_x\right)
       H_{\rm marks}(w),
\end{equation}
where $\psi_1$ has fixed compact support and its rescaled derivatives,
as well as $|\sigma|$, have fixed log-power bounds.  For every fixed
$B'>0$, we claim uniformly for $|t|\leq L^{B'}$ that the normalized
sum in \eqref{cor:SM} saves any prescribed log power.

In a tuple term put $w=ln$.  All long primes lie outside the marked
groups, so $H_{\rm marks}(ln)=H_{\rm marks}(n)$.  Apply
Lemma~\ref{lem:divisor-model} on the scales $Y/l$ and $Y$ in the
tuple and constant terms respectively.  Cauchy and
\eqref{cor:model-error} make the total normalized error at most
$L^{-N}(1+\sum_{\mathbf l}1/l)$, after renaming $N$.
For a fixed divisor $d'$, all primes in $ld'$ exceed $k$ for large
$x$, so $\gcd(ld',k)=1$.  Smooth sum--integral comparison on the
progressions modulo $k$ gives
\begin{align}
 &\frac1Y\sum_{v\geq1}
   \psi_1(ld'v/Y)(ld'v)^{i(t+\sigma)}\e(hld'v/k)
 \notag\\
 &\quad=\frac1{ld'}
       \left(\frac1k\sum_{a\bmod k}\e(ha/k)\right)
       \int_0^\infty\psi_1(z)(Yz)^{i(t+\sigma)}\,\dd z
       +O(Y^{-1}L^{C(B')}).
 \label{cor:low-main-term}
\end{align}
The sum has scale $Y/(ld')\geq x^{\eps-o(1)}$.  The error follows
by summing the endpoint and total-variation errors over $k$ residue
classes; all derivative and modulus costs are fixed log powers.
Moreover
\[
 \sum_{\mathbf l}1\leq x^{1-\eps}C_x,\qquad
 \sum_{d'}|A_{d'}|
 \leq\exp(O(L^{.47}))L^{C_8}=x^{o(1)}.
\]
The summed errors in \eqref{cor:low-main-term} are therefore
$O(x^{-\eps+o(1)})$, since $Y\geq x^{1-o(1)}$.
The main term depends on $l$ only through $1/l$.  Its tuple sum is
exactly $C_x$ times the main term for the subtracted constant.  They
cancel, including the complete residue average.  Choosing $N$ after
$B'$ proves \eqref{cor:SM} on every fixed log-power time interval.

\subsection{High frequencies: factorization and exceptional times}

Choose a small group $g_s$.  It has $t_{g_s}=1$ in every comparison,
and the tuple test depends only on big marks.  Except when a prime of
this group divides $w$ twice,
\begin{equation}\label{cor:extract-small}
 W_{\mathbf t}^{\mathbf b}(w)
  =\frac1{V_{g_s}}\sum_{\substack{p_s\mid w\\p_s\in\mathcal P_{g_s}}}
       W_{\mathbf t-\mathbf e_{g_s}}^{\mathbf b}(w/p_s).
\end{equation}
This follows directly by removing the represented small prime; the
remaining damping exponent is unchanged.  Both sides are bounded by
fixed log powers.  The square exceptions have relative count at most
$L^{O(1)}\exp(-L^{.27})$ on the relevant dyads.  In the positive
tuple part we may also allow repeated slot primes: the extra terms
have a square of a prime at least $x^\eps$ dividing $w$, a relative
count $O(x^{-\eps})$.  Slot multiplicities at each $w\asymp Y$ remain
bounded.  Lemma~\ref{lem:moment-bounds}, applied to these coefficient
errors, makes their normalized integrated square over $|t|\leq T_0$
smaller than every negative log power.  Indeed their squared coefficient
norm is at most
$YL^{O(1)}(\exp(-L^{.27})+x^{-\eps})$, and
$T_0/Y=L/H=o(1)$.

In the tuple part factor $w=p_sp_lu$, where $p_l=l_1$ lies in the
first long-prime slot.  After allowing repeated slots the coefficient
on $u$ is exactly
\begin{equation}\label{cor:u-coefficient}
 W_{\mathbf t-\mathbf e_{g_s}}^{\mathbf b}(u)\mathcal W(u)
       \sum_{l_2,\ldots,l_d}\1_{l_2\cdots l_d\mid u}.
\end{equation}
For $d=1$ the final factor is one.  This is bounded by a fixed log
power times a fixed divisor power.  In the constant part write
$w=p_sn$; its coefficient is just the remaining marked weight times
$\mathcal W(n)$.

Separate the rational phase by residue classes modulo $k$.  In the
tuple part, $p_s,p_l$ are units modulo $k$; on fixed classes for $p_s$
and $u$, expand the remaining test in characters of $p_l$ by
orthogonality on the units.  In the constant part fix classes for
$p_s,n$.  Both operations have polynomial-in-$k$ total cost and do
not require $u$ or $n$ to be units.  Divide all factors into dyads and
Fourier-separate the smooth product cutoff in logarithmic size.
There are only a fixed log-power number of boxes.  Their factor scales
have product comparable to $Y$, and every collapsed coefficient
sequence of normalized scale $R$ has
\begin{equation}\label{cor:collapsed-coefficients}
 \sum_{v\asymp R}|c_v|^2\ll L^{C_{11}}/R,
\end{equation}
by fixed divisor moments.  This applies also to the uncut product of
all factors, so Dirichlet mean squares uniformly on translated time
intervals justify discarding the Fourier tails with arbitrary log
saving.  Retain only shifts of fixed log-power size.  Enlarge the
low-time exponent $B'$ beyond these shifts and any later prime
thresholds.  The low-time argument already proved is available for
this enlarged $B'$.  At the remaining times all retained shifts of $t$
are comparable in magnitude to $t$.

It remains to treat products of normalized polynomials at a common
time, with $L^{B'}\leq|t|\leq2T_0$.  In every such product there is
a small-prime factor
\begin{equation}\label{cor:small-polynomial}
 P_s(t)=P^{-1}\sum_{\substack{p_s\asymp P\\p_s\in\mathcal P_{g_s}}}
            b_s(p_s)p_s^{it},\qquad
 cL^{.27}\leq\log P\leq L^{.36},\qquad |b_s|\leq L^{C_{12}}.
\end{equation}
The coefficients include residue restrictions and retained twists.
In the tuple part the remaining factors are a long-prime polynomial
$P_l(t)$ and a polynomial $R(t)$ at scales $N_l$ and
\begin{equation}\label{cor:remaining-scale}
 R_0\asymp Y/(PN_l)\geq x^{\eps-o(1)}.
\end{equation}
In the constant part the remaining polynomial $N(t)$ has scale
$R_0\asymp Y/P$.

The threshold split and exceptional-frequency treatment follow the
Matom\"aki--Radziwi{\l}{\l} strategy; see
\cite[\S2.1 and Lemmas~8--9 of arXiv~v4]{MatomakiRadziwill2017}.
The residual-scale sparse Gram estimate below is adapted here using
Appendix~A, while the factorial coefficient estimate retains the
Soundararajan credit stated at its use.

On the set $|P_s(t)|\leq L^{-A_s}$, collapse all remaining factors.
Their scale is $Y/P$, and
\begin{equation}\label{cor:length-dominates-time}
 \frac{Y/P}{2T_0}=\frac H{2LP}\longrightarrow\infty,
\end{equation}
because $H\geq\exp(2L^{.39}-O(\log L))$ and
$P\leq\exp(L^{.36})$.  Equations
\eqref{cor:collapsed-coefficients} and the mean-square bound give a
fixed log-power integrated square for this collapsed polynomial.
Taking $A_s$ sufficiently large handles this part.

Let $\mathcal J$ be the integer unit intervals meeting
$\{|t|\leq2T_0:|P_s(t)|>L^{-A_s}\}$.  Then
\begin{equation}\label{cor:exceptional-count}
 |\mathcal J|\leq
 \exp\bigl(O(L^{.36}+L^{.73}\log L)\bigr)
       \leq\exp(O(L^{.9})).
\end{equation}
For this step we use the factorial prime-polynomial coefficient
estimate, as in \cite[Lemma~3 and its proof]{Soundararajan2009},
together with a coefficient-independent mean square. Put
$Z=8T_0+8$ and $j_0=\lfloor\log Z/\log(2P)\rfloor$.  Unique
factorization gives squared coefficient norm for $P_s^{j_0}$ at most
\[
 j_0!\left(P^{-2}\sum_{p_s\asymp P}|b_s(p_s)|^2\right)^{j_0}
       \leq j_0!(L^{C_{13}}/P)^{j_0}.
\]
Select a point exceeding the threshold from each occupied interval,
and split their integer indices modulo three to get separated sets.
The indices of this powered polynomial are at most $Z$.  The
separated-point mean-square bound of Lemma~\ref{lem:moment-bounds}
therefore gives
\[
 |\mathcal J|\ll ZL^{O(1)}j_0!
                       (L^{C_{13}+2A_s}/P)^{j_0}.
\]
Since $ZP^{-j_0}\leq2P\,2^{j_0}$ and
$j_0=O(L^{.73})$, taking logarithms proves
\eqref{cor:exceptional-count}.

We will also use the following sparse bound for any normalized
polynomial $R(t)=\sum_{n\asymp R_0}c_nn^{it}$ with
$x^{\eps/2}\leq R_0\leq x^2$ and
$\sum|c_n|^2\leq L^{C_{14}}/R_0$:
\begin{equation}\label{cor:sparse-bound}
 \sum_{I\in\mathcal J}\sup_{t\in I}|R(t)|^2\ll L^{C_{15}}.
\end{equation}
Choose maximizing points and again split into three separated classes.
For $1\leq|\Delta|\leq T_{\rm mod}:=\exp(L/(\log L)^2)$,
sum--integral comparison on the full integer dyad gives
\[
 R_0^{-1}\left|\sum_{n\asymp R_0}n^{i\Delta}\right|
       \ll |\Delta|^{-1}+x^{-c_\eps}.
\]
For larger differences, Lemma~\ref{lem:log-phase} gives
$O(\exp(-L/(\log L)^{C_{16}}))$ instead.  Its hypotheses hold:
$R_0\geq x^{\eps/2}>\exp(L/(\log L)^2)$, all relevant scales are
at most $2x^5$, and the differences are between $T_{\rm mod}$ and
$4T_0+O(1)<4x^3$.  The Gram matrix of the vectors
$(n^{it_i})_{n\asymp R_0}$ consequently has absolute row sum
\[
 \ll R_0\left(1+\log(2T_0)
       +|\mathcal J|\{x^{-c_\eps}
                   +\exp(-L/(\log L)^{C_{16}})\}\right)
 \ll R_0L.
\]
Here separation bounds the reciprocal-difference sum by $O(\log T_0)$,
and \eqref{cor:exceptional-count} absorbs both small errors.  The Gram
operator bound multiplied by $\sum|c_n|^2$ proves
\eqref{cor:sparse-bound}.

\subsection{The tuple and constant contributions at high times}

In the tuple term the long-prime polynomial is, after partial summation,
a normalized prime sum with a character modulo $k$ and a fixed
log-power twist.  Its dyad lies between $x^\eps/2$ and
$x^{1-\eps}$.  Choose fixed $0<\tau<\eps$ and
$1-\eps<\eta<1$ to apply Lemma~\ref{lem:long-prime}.  For any
prescribed $A_l$, it gives
\begin{equation}\label{cor:long-prime-saving}
 |P_l(t)|\ll L^{-A_l}
       \quad(L^{B'}\leq|t|\leq2T_0),
\end{equation}
after increasing $B'$ beyond its threshold and all retained shifts.
The upper frequencies are $<x^2$.  The remaining coefficient
\eqref{cor:u-coefficient} has scale \eqref{cor:remaining-scale} and
obeys \eqref{cor:collapsed-coefficients}, so
\eqref{cor:sparse-bound} applies.  The small polynomial has absolute
bound $L^{O(1)}$.  Hence on the exceptional intervals the integrated
square of $P_sP_lR$ is
\[
 \ll L^{O(1)-2A_l}
       \sum_{I\in\mathcal J}\sup_{t\in I}|R(t)|^2,
\]
which has arbitrary log saving on choosing $A_l$ large enough.
Together with the ordinary-time estimate this handles the tuple part.

For the constant part set $R_0\asymp Y/P$, and apply
Lemma~\ref{lem:divisor-model} to the remaining marked coefficient at
this scale.  Keep its residue-class restriction separately.  If $E(n)$
is the coefficient error, then for every fixed $N$,
$\sum_{n\asymp R_0}|E(n)|^2\ll R_0L^{-N}$.  The mean-square bound
on the \emph{entire} retained time interval gives
\begin{align}
 \int_{|t|\leq2T_0}
   \left|R_0^{-1}\sum_{n\asymp R_0}E(n)n^{it}\right|^2\,\dd t
 &\ll L^{-N}\left(\frac{T_0}{R_0}+\log R_0\right)\notag\\
 &\ll L^{-N}\left(\frac{LP}{H}+O(L)\right).
 \label{cor:constant-model-error}
\end{align}
By \eqref{cor:length-dominates-time}, $LP/H=o(1)$.  Thus the known
small-prime suprema and all fixed decomposition costs can be paid by
choosing $N$ last.  In particular the model error is controlled before
restricting to the exceptional times.

For the truncated model, a divisor $d'$ leaves a progression modulo
$k$ at scale $R'=R_0/d'$.  Its normalization in $N(t)$ is
$A_{d'}/d'$ times normalization at $R'$.  Since
$d'\leq\exp(O(L^{.47}))$, one has $R'=x^{1-o(1)}$.  The divisor
is a unit modulo $k$, so the original residue restriction becomes a
single residue restriction at this scale.  For
$L^{B'}\leq|t|\leq T_{\rm mod}$, comparison with the integral of
the logarithmic phase gives
\begin{equation}\label{cor:constant-medium}
 |N_0(t)|\ll L^{C_{17}}\bigl(|t|^{-1}+x^{-c'}\bigr).
\end{equation}
Indeed the normalized integral over a dyad is $O(1/|t|)$ and the
progression endpoint and variation error is
$O(L^{O(1)}(1+|t|)/R')=O(x^{-c'})$; sum using
$\sum|A_{d'}|/d'\leq L^{C_8}$.  Retained twists may either be
incorporated in $t$ or absorbed by increasing $B'$.
The square of \eqref{cor:constant-medium} is integrable with
\begin{equation}\label{cor:constant-medium-integral}
 \int_{L^{B'}\leq|t|\leq T_{\rm mod}}|N_0(t)|^2\,\dd t
       \ll L^{2C_{17}-B'}+x^{-c''}.
\end{equation}
This has arbitrary log saving after enlarging $B'$, even after the
small-prime absolute bound.  We used square integration of $1/|t|$,
not a pointwise log saving multiplied by the full interval length.

For $T_{\rm mod}<|t|\leq2T_0$, apply
Lemma~\ref{lem:log-phase} to each such progression.  Its modulus is
at most $L^{C_6}$, its scale is $x^{1-o(1)}$, and, including retained
shifts, its frequency is at least
$\tfrac12T_{\rm mod}>\exp(L/(2(\log L)^2))$ and at most
$x^{1+o(1)}<4x^3$.
After summing the reciprocal divisor coefficients this yields
\begin{equation}\label{cor:constant-high}
 |N_0(t)|\ll L^{C_{18}}\exp(-L/(\log L)^{C_{19}}).
\end{equation}
We now integrate only on the exceptional unit intervals.  By
\eqref{cor:exceptional-count} and the small-prime absolute bound,
their total contribution is at most
\[
 L^{O(1)}\exp\bigl(O(L^{.9})-2L/(\log L)^{C_{19}}\bigr),
\]
which saves every fixed log power.  No length-dominated mean-square
estimate was required on an individual divisor progression: $R'$ may
be smaller than $T_0$, and the bounds
\eqref{cor:constant-medium}--\eqref{cor:constant-high} are pointwise.
Equations \eqref{cor:constant-model-error},
\eqref{cor:constant-medium-integral}, and \eqref{cor:constant-high}
complete the constant-term estimate.

We have proved \eqref{cor:SM} to every fixed log precision.  The local
Mellin inequality gives \eqref{cor:arc-energy-target}; the comparison
multiplier estimate and Parseval then bound each term in
\eqref{cor:comparison-sum}, divided by $Y$, by an arbitrarily large
negative log power.  Choose that precision after the known fixed
kernel, dyad, arc, and separation costs.  Their total is therefore
negligible.  Subtracting these comparison terms from the residual
pairing leaves its raw pattern, already identified with
\eqref{eq:two-sided}.  This proves Theorem~\ref{thm:two-sided}.
